\par
% Certificate constants.
% name                         value / source
\newcommand{\brM}{77}
\newcommand{\brModes}{72}
\newcommand{\brWindow}{160}
\newcommand{\brStates}{300}
\newcommand{\brGaussNodes}{1536}
\newcommand{\brDGrid}{512}
\newcommand{\brTraceGrid}{1024}
\newcommand{\brTraceCutoff}{400}
\newcommand{\brSpectralBits}{2048}
\newcommand{\brTraceBits}{4096}
\newcommand{\brRecordedSeconds}{173.042}
\newcommand{\brWorkers}{4}
\newcommand{\brCentreHash}{272abaf4b002d8b63cf68344616ca501eb1a3c1f2636b27280af4960713c4472}
\newcommand{\brVerifierHash}{375ea2e68225d3419309d5359c7d5d2a589c07bb0985e71f2cd0ac9bc3fc95fb}
\newcommand{\brRadius}{10^{-36}}
\newcommand{\brRhoLo}{237.97}
\newcommand{\brRhoHi}{237.98}
\newcommand{\brAnalyticRadius}{1.00904}
\newcommand{\brDerivativeLo}{237.2785}
\newcommand{\brUnivalentLo}{235.7747}
\newcommand{\brConvexLo}{0.64407}
\newcommand{\brNonDiscLo}{0.0038627}
\newcommand{\brOriginLo}{19.3341}
\newcommand{\brMu}{0.011202}
\newcommand{\brEpsC}{5.276\cdot10^{-13}}
\newcommand{\brEps}{2.580\cdot10^{-11}}
\newcommand{\brSchur}{0.013407}
\newcommand{\brDinf}{0.355923}
\newcommand{\brDiagLo}{0.5561235}
\newcommand{\brGaussError}{9.734\cdot10^{-51}}
\newcommand{\brBLtwo}{1.823}
\newcommand{\brBY}{1.117\cdot10^6}
\newcommand{\brBTone}{6.443\cdot10^{10}}
\newcommand{\brBTtwo}{1.187\cdot10^{13}}
\newcommand{\brHighFactor}{0.033901}
\newcommand{\brP}{240.1776}
\newcommand{\brPone}{295.0400}
\newcommand{\brPtwo}{42522.831}
\newcommand{\brJtwo}{0.793080}
\newcommand{\brBzeroone}{541.556}
\newcommand{\brUb}{68.459}
\newcommand{\brDUb}{17609.075}
\newcommand{\brEzero}{9.699\cdot10^{-67}}
\newcommand{\brEone}{2.728\cdot10^{-62}}
\newcommand{\brA}{1.947\cdot10^{22}}
\newcommand{\brYzero}{1.888\cdot10^{-44}}
\newcommand{\brZ}{2.981\cdot10^{-39}}
\newcommand{\brMtwo}{65827.163}
\newcommand{\brCtwo}{6.406\cdot10^{26}}
\newcommand{\brPolynomial}{-9.99999980\cdot10^{-37}}
\newcommand{\brContraction}{1.282\cdot10^{-9}}
\newcommand{\brFourCY}{4.836\cdot10^{-17}}
\newcommand{\brCoverageRows}{%
Angular order & $0$ & $77$ & $154$ & $231$ & $308$ & $385$ & $462$\\
\midrule
States & $102$ & $91$ & $60$ & $33$ & $13$ & $1$ & $0$\\}
% Floating-point diagnostics.
\newcommand{\brCurvatureDiagnostic}{[2.719\cdot10^{-3},\,5.283\cdot10^{-3}]}
\newcommand{\brCSSConvexDiagnostic}{-2.5506}
\newcommand{\brCSSCurvatureDiagnostic}{-2.968}
\newcommand{\brAlphaDiagnostic}{56632.65}
\newcommand{\brScaleDiagnostic}{84.137}
\newcommand{\brRemainderDiagnostic}{8467}
\newcommand{\brThresholdDiagnostic}{92.6\text{--}228.8}
\newcommand{\brRatioDiagnostic}{37}
\newcommand{\brRemainderRows}{%
$5000$ & $3953$\\
$23309.55191881634$ & $7096$\\
$100000$ & $8467$\\}
% END BERENSTEIN CERTIFICATE MACRO TABLE

\part{A convex planar domain for Berenstein's problem}\label{br:part}
\section{The domain and the fixed-disc equation}\label{br:sec:formulation}

The Dirichlet counterpart of Schiffer's problem prescribes zero boundary
value and a nonzero constant normal derivative. Colbrook, Sadeghi and
Stepaniants~\cite{CSS26} constructed a planar counterexample to
Berenstein's conjecture. The construction below gives the first convex
planar counterexample.

\begin{theorem}\label{br:thm:main}
There exist a bounded strictly convex domain $\Omega\subset\R^2$, different
from a disc, with real-analytic boundary and $D_{\brM}$ symmetry, and a
real-valued sign-changing function $u\in C^\infty(\overline\Omega)$ such that
\begin{equation}\label{br:eq:physical}
 \Delta u+u=0\quad\text{in }\Omega,\qquad
 u=0,\quad\partial_\nu u=1\quad\text{on }\partial\Omega.
\end{equation}
Here $\nu$ is the outward unit normal.
The domain is $\psi(\D)$, with $\psi(0)=0$, $\psi'(0)>0$, and
\begin{equation}\label{br:eq:main-geometry}
 \brRhoLo<\psi'(0)<\brRhoHi,\qquad
 \re\left(1+\frac{w\psi''(w)}{\psi'(w)}\right)>\brConvexLo
 \quad(|w|=1).
\end{equation}
The map is holomorphic and univalent on $|w|<R=2^{1/\brM}>\brAnalyticRadius$,
with $\re\psi'(w)>\brUnivalentLo$ there. On $\overline\D$,
$|\psi'|>\brDerivativeLo$. Writing
$\psi(w)=\sum_{j\ge0}c_jw^{jm+1}$, $m=\brM$, one has
$|c_1|>\brNonDiscLo$ and $u(0)>\brOriginLo$.
There is a unique zero of the fixed-disc equation~\eqref{br:eq:F} in
\begin{equation}\label{br:eq:solution-ball}
 \nrm{g-g_0}_Y+\nrm{p-p_0}_{W_1}\le\brRadius,
 \qquad p=\psi'.
\end{equation}
Here $Y$, $W_1$ and $X$ are the spaces of
Subsection~\ref{br:sec:spaces}, and $(g_0,p_0)$ is the clamped centre of
Subsection~\ref{br:sec:centre}.
\end{theorem}

The second quantity in~\eqref{br:eq:main-geometry} is the conformal convexity
factor. The physical curvature is this factor divided by $|\psi'|$.
For the finite centre, a floating-point evaluation gives physical curvature
in the approximate range $\brCurvatureDiagnostic$.

\subsection{The arc-length Fourier transform}\label{br:sec:fourier}
Let $\sigma_{\partial\Omega}$ denote arc-length measure and use the convention
$\widehat\sigma(\xi)=\int_{\partial\Omega}e^{-i\xi\cdot x}\,ds$.
For a smooth bounded
domain with connected Jordan boundary, the existence of a real solution
$u$ of~\eqref{br:eq:physical} is equivalent to
\begin{equation}\label{br:eq:fourier}
 \widehat\sigma_{\partial\Omega}(\xi)=0\qquad(|\xi|=1).
\end{equation}
Indeed, if $v=e^{-i\xi\cdot x}$, Green's identity gives
$0=\int_\Omega(u\Delta v-v\Delta u)
=\int_{\partial\Omega}(u\partial_\nu v-v\partial_\nu u)
=-\widehat\sigma(\xi)$. Conversely, let $\Phi$ be the outgoing fundamental
solution with $(\Delta+1)\Phi=-\delta$ and put
$S(x)=\int_{\partial\Omega}\Phi(x-y)\,ds_y$.
Its far-field amplitude is a nonzero constant multiple of
$\widehat\sigma(\widehat x)$, so~\eqref{br:eq:fourier} and Rellich uniqueness
give $S=0$ in the connected exterior. The single layer is continuous across
the boundary, and its normal derivatives obey
$\partial_\nu S|_{\rm ext}-\partial_\nu S|_{\rm int}=-1$.
Thus $u=\re S|_\Omega$ solves~\eqref{br:eq:physical}; smooth boundary
regularity gives the classical traces. This is the Fourier-circle
equivalence of~\cite[Proposition~2.1]{CSS26}. In particular, the theorem produces a strictly
convex noncircular curve whose arc-length Fourier transform vanishes on
the unit circle.

\subsection{Spaces and lifts}\label{br:sec:spaces}
In this part $m=\brM$, $R=2^{1/m}$, and the symbols $Y$, $\Yc$, $K$, $X_g$
have the meanings of Subsection~\ref{r2:sec:spaces}, with this value of $m$.
The field is real, reflection-even and invariant under rotation through
$2\pi/m$. In particular,
\begin{equation}\label{br:eq:norms}
 \nrm f_Y=\sum_{n\in m\Z}R^{|n|}\sup_{0\le r\le1}|f_n(r)|,
 \qquad
 \nrm b_{W_s}=\sum_{n\in m\Z}R^{|n|}(1+|n|)^s|b_n|.
\end{equation}
The full angular space $\Yc$ permits all integer frequencies. Continuity
is Cartesian, so nonconstant angular coefficients vanish at the origin.
Holomorphic versions of $W_s$ use nonnegative frequencies and real
coefficients. The algebra proof of Lemma~\ref{r2:lem:algebra} applies for
every $R>1$ and every $m$; the additional weight satisfies
$1+|n+k|\le(1+|n|)(1+|k|)$, making $W_s$ a Banach algebra for
$s=0,1,2$. A holomorphic real-coefficient series $h$ satisfies
\begin{equation}\label{br:eq:realpart-norm}
 \nrm{\re h}_{W_s}=\nrm h_{W_s}.
\end{equation}
For each positive frequency, taking the real part puts half the Taylor
coefficient at each of the two opposite Fourier frequencies; their weights
are equal. The constant coefficient is real and retains its weight.

The zero-Dirichlet Poisson inverse $K$ and its clamped range are
\begin{equation}\label{br:eq:clamped-space}
 X_g=\left\{g\in Y:\int_0^1r^{|n|+1}g_n(r)\,dr=0\ (n\in m\Z)\right\},
 \qquad \nrm K_{Y\to Y}\le\tfrac14.
\end{equation}
Lemma~\ref{r2:lem:Poisson} uses neither the Schiffer boundary value nor the
symmetry order: its radial Green formula and moment identity apply
verbatim. Thus $Kg$ has zero Cauchy data precisely when $g\in X_g$.
For a boundary series introduce the Neumann lift
\begin{equation}\label{br:eq:neumann-lift}
 \Bc b=\frac12\sum_n b_nr^{|n|}(r^2-1)e^{in\theta}.
\end{equation}
Its Cauchy traces are $(0,b)$, and
\begin{equation}\label{br:eq:neumann-lift-bounds}
 \nrm{\Bc b}_Y\le\tfrac12\nrm b_{W_0},\qquad
 \nrm{\Delta\Bc b}_Y\le2\nrm b_{W_1}.
\end{equation}
Indeed $\Delta(r^{|n|}(r^2-1)e^{in\theta}/2)
=2(|n|+1)r^{|n|}e^{in\theta}$, and the radial factor has modulus at most
$1/2$ on $[0,1]$.

Write $p=\psi'=\sum_jp_jw^{jm}$ and
$\psi(w)=\int_0^wp(z)\,dz$, so $c_j=p_j/(jm+1)$.
Conformal covariance~\eqref{r2:eq:covariance} and the normal transformation
give the following fixed-disc problem for $U=u\circ\psi$:
\begin{equation}\label{br:eq:pullback}
 \Delta U+|p|^2U=0\quad\text{in }\D,\qquad
 U=0,\quad U_r=|p|\quad\text{on }\partial\D.
\end{equation}
The boundary vector $wp(w)$ is outward normal to the image and has length
$|p(w)|$, so the normal derivative in physical coordinates is
$|p|^{-1}U_r$.
We work on the real Banach space
\begin{equation}\label{br:eq:X}
 X=X_g\times W_1,\qquad
 \nrm{(g,p)}_X=\nrm g_Y+\nrm p_{W_1},\qquad
 U(g,p)=\Bc(|p|)+Kg.
\end{equation}
The shape component $p$ of $X$ is a holomorphic real-coefficient series
in $W_1$: $p=\sum_{j\ge0}p_jw^{jm}$ with $p_j\in\R$. Its reflection
symmetry $p(\bar w)=\overline{p(w)}$ is equivariance; it does not make
$p$ real on $\partial\D$. The field is real-valued and reflection-even
in $Y$. On the ball~\eqref{br:eq:solution-ball}, the boundary function
$|p|$ lies in the full two-sided space $W_1$, and every shape satisfies
$\re p\ge2\rho-P-\brRadius>0$ on the closed weighted disc $|w|\le R$,
with $P$ of~\eqref{br:eq:P}; see~\eqref{br:eq:univalence}. Thus $p$ is
nonzero throughout that disc.
We solve
\begin{equation}\label{br:eq:F}
 F(g,p)=g+\Delta\Bc(|p|)+|p|^2\bigl(\Bc(|p|)+Kg\bigr)=0.
\end{equation}
On the ball $\nrm{p-p_0}_{W_1}<1/J_1$, with $J_1$ as
in~\eqref{br:eq:reciprocals}, write $p=p_0(1+k/p_0)$. The binomial
series for the square roots of $1+k/p_0$ and its reflection converge
in $W_1$, and $|p_0|\in W_1$ by~\eqref{br:eq:b-bound}. Thus $F$
is a real-analytic map from this neighbourhood of the centre in $X$ to $Y$.
The term $\Delta\Bc(|p|)$ requires the $W_1$ shape norm.
For every $(g,p)\in X$ in this neighbourhood, $U(g,p)$ has the exact
Cauchy data $(0,|p|)$, and $F=\Delta U+|p|^2U$ distributionally.
The Schiffer cubic equation and its Dirichlet linear inverse are therefore
replaced by~\eqref{br:eq:F} and a Robin inverse.

\subsection{The exactly clamped centre}\label{br:sec:centre}
The input consists of exact decimal coefficients
$c_j^{(0)}$, $a_j$ and positive divisors $s_j$, $0\le j\le\brModes$.
The divisors are specified positive numbers; no identity with a Bessel
value is assumed. The decimal strings are the entries $0\le j\le72$
of \texttt{p}, \texttt{a} and \texttt{scales} in
\path{anc/berenstein/centre_frozen.json}; the array \texttt{p} stores
$c_j^{(0)}$, while $p_j^{(0)}=(jm+1)c_j^{(0)}$. Put
\begin{equation}\label{br:eq:input}
 \begin{split}
 \psi_0(w)&=\sum_{j=0}^{\brModes}c_j^{(0)}w^{jm+1},\quad
 p_0=\psi_0',\quad \rho=c_0^{(0)},\quad q=|p_0|^2,\\
 u_*(z)&=\sum_{j=0}^{\brModes}\frac{a_j}{s_j}
 J_{jm}(|z|)\cos(jm\arg z),\quad U_*=u_*\circ\psi_0.
 \end{split}
\end{equation}
Each physical summand extends through the origin as an entire real-analytic
Helmholtz solution, so $(\Delta+q)U_*=0$ identically.
The finite conformal polynomial and finite field specify an infinite-dimensional
clamped centre as follows. Let $D_j$ and $N_j$ be the complete cosine
coefficients of
\begin{equation}\label{br:eq:defects}
 U_*(1,\theta)=\sum_{j\ge0}D_j\cos(jm\theta),\qquad
 \partial_rU_*(1,\theta)-|p_0(e^{i\theta})|
 =\sum_{j\ge0}N_j\cos(jm\theta).
\end{equation}
For $n\ge0$ use the polynomials of~\eqref{r2:eq:lifts},
\[
 h_n^D=r^n\bigl(1-n(r^2-1)/2\bigr),\qquad
 h_n^N=r^n(r^2-1)/2,
\]
whose traces $(1,0)$ and $(0,1)$ do not depend on the prescribed Cauchy
data. Define
\begin{equation}\label{br:eq:centre}
 \begin{split}
 \ell&=\sum_{j\ge0}(D_jh_{jm}^D+N_jh_{jm}^N)\cos(jm\theta),\\
 U_0&=U_*-\ell,\qquad
 g_0=\Delta\bigl(U_0-\Bc(|p_0|)\bigr).
 \end{split}
\end{equation}
The strip estimates of Subsection~\ref{br:sec:residual} give convergence with
all derivatives needed near $\overline\D$. Thus $U_0$ has traces
$(0,|p_0|)$, the quantity inside the Laplacian has zero Cauchy data,
$g_0\in X_g$, and Poisson uniqueness gives
$Kg_0=U_0-\Bc(|p_0|)$. In particular $U(g_0,p_0)=U_0$ and
\begin{equation}\label{br:eq:E}
 E:=F(g_0,p_0)=-(\Delta+q)\ell.
\end{equation}
All boundary harmonics enter this definition. The exact zero enclosed
in~\eqref{br:eq:solution-ball} need not have a finite conformal or field
series.
\section{The Robin inverse in the symmetric sector}\label{br:sec:spectrum}
At the centre define
\begin{equation}\label{br:eq:Robin}
 LV=(\Delta+q)V=f,\qquad V_r+\beta V=0\quad\text{on }\partial\D,
 \qquad \beta=1+d,\quad d=\re\frac{wp_0'}{p_0}.
\end{equation}
All operators in this section act in the real even $D_m$ sector, with
interior measure $r\,dr\,d\theta/(2\pi)$ and boundary measure
$d\theta/(2\pi)$. By~\eqref{r2:eq:curvature-formula}, $\beta=|p_0|\kappa_0$,
where $\kappa_0$ is the curvature of $\partial\psi_0(\D)$. The symmetric form associated with $-L$ is
\begin{equation}\label{br:eq:Robin-form}
 \mathfrak a(z,y)=\int_\D\nabla z\cdot\nabla y
 +\int_{\partial\D}\beta zy-\int_\D qzy,
 \qquad z,y\in H^1(\D).
\end{equation}
The form is closed and bounded below after addition of a sufficiently large
multiple of the $L^2$ inner product. Its self-adjoint operator is the Robin
realisation of $-L$. Boundary multiplication is treated as a form on $H^1$.

\subsection{The Dini basis and complete spectral window}\label{br:sec:coverage}
Let $\alpha_{j,k}>0$ run over the roots
\begin{equation}\label{br:eq:Dini-roots}
 \alpha J_n'(\alpha)+J_n(\alpha)=0,\qquad n=jm.
\end{equation}
The Robin-one Laplacian $-\Delta_{R,1}$ has the complete orthonormal Dini
basis
\begin{equation}\label{br:eq:basis}
 e_{j,k}=b_{j,k}(r)\chi_j(\theta),\qquad
 b_{j,k}(r)=\frac{\sqrt2J_n(\alpha_{j,k}r)}
 {|J_n(\alpha_{j,k})|\sqrt{1+(1-n^2)/\alpha_{j,k}^2}},
\end{equation}
where $\chi_0=1$ and $\chi_j=\sqrt2\cos(jm\theta)$ for $j>0$.
The Dini parameter is $1$, with $n+1>0$, so no exceptional nonpositive
radial eigenvalue is present. Orthogonality and completeness also follow
from the separated self-adjoint Robin Sturm--Liouville problems and the
complete angular cosine basis. The identity
\[
 \int_0^1rJ_n(\alpha r)^2\,dr
 =\tfrac12\bigl(J_n'(\alpha)^2+(1-n^2/\alpha^2)J_n(\alpha)^2\bigr)
\]
follows by integrating the Bessel equation and its parameter derivative.
At~\eqref{br:eq:Dini-roots} it gives
$\tfrac12J_n(\alpha)^2(1+(1-n^2)/\alpha^2)$, which is positive:
every root exceeds $n$ for $n\ge1$, and the $n=0$ expression is positive.
In particular the basis functions have nonzero boundary values.

Let $\PP$ select every root satisfying
\begin{equation}\label{br:eq:window}
 |\alpha_{j,k}-\rho|<\omega,\qquad \omega=\brWindow,
\end{equation}
and put $\QQ=I-\PP$. The window contains $\brStates$ states, with the
sector counts of Table~\ref{br:tab:coverage}.
Enumerate these states by $i\leftrightarrow(j,k)$ and write
$n_i=jm$, $\alpha_i=\alpha_{j,k}$, $\chi_i=\chi_j$ and
$\tau_i=b_{j,k}(1)$.
\begin{table}[htbp]
\centering
\begin{tabular}{c rrrrrrr}
\toprule
\brCoverageRows
\bottomrule
\end{tabular}
\caption{Complete Robin-one spectral window. All higher angular orders
have no state in the window.}\label{br:tab:coverage}
\end{table}

\begin{lemma}[Two coverage arguments]\label{br:lem:coverage}
The Dirichlet interlacing and the quarter-step Robin sign grid described
below enumerate all states of~\eqref{br:eq:window}.
\end{lemma}
\begin{proof}
For $y(x)=\sqrt xJ_n(x)$ the Bessel equation gives
\[
 y''+Q_n(x)y=0,\qquad Q_n(x)=1-(n^2-1/4)/x^2.
\]
On an interval with $Q_n\le Q_{\max}$, integration by parts between two
zeros and Wirtinger's inequality give their spacing at least
$\pi/\sqrt{Q_{\max}}$, exactly as in
Lemma~\ref{r2:lem:zero-coverage}. Enclose proposed Dirichlet zeros by
opposite signs in disjoint brackets of width less than $1/10$.
For successive brackets the distance between their outer endpoints is
less than $2\pi/\sqrt{Q_{\max}}$, where $Q_{\max}$ is the maximum
between those endpoints. Three zeros there would require the reverse
inequality, so each bracket contains exactly one zero and none lies
between them. Guard zeros beyond both endpoints of the window
complete this argument whenever a lower Dirichlet zero is present.

For every positive angular order whose first Dirichlet zero lies above
the lower window endpoint, the initial interval requires a separate guard.
This includes $n=m,2m,\ldots,5m$ in the present window. When $n>0$, $J_n$ and $J_n'$ are
positive on $(0,n]$: at a first derivative zero below $n$ the Bessel ODE
would give $J_n''=(n^2/x^2-1)J_n>0$, which cannot be a first downward
crossing. This also excludes a first zero of $J_n$ there. The identity
$(xJ_n')'=(n^2/x-x)J_n>0$ on $(0,n)$ gives $J_n'(n)>0$,
and hence $y(n),y'(n)>0$. Constant-coefficient comparison puts the
first zero after $n$ at distance at least $\pi/(2\sqrt{Q_{\max}})$.
A hidden first zero before the proposed bracket would put the bracketed
zero at distance at least $3\pi/(2\sqrt{Q_{\max}})$. The strict opposite
inequality is imposed on its upper endpoint, with $Q_{\max}$ taken on
the entire interval from $n$ to that endpoint. This proves the initial
Dirichlet guard.

Between consecutive Dirichlet zeros consider $z=xJ_n'/J_n$. It satisfies
$xz'=n^2-x^2-z^2$. Thus $1+z$ runs from $+\infty$ to $-\infty$ and, at any
zero, has derivative $-(x^2+1-n^2)/x<0$. There is exactly one Robin root
in each such gap. In the initial gap it is likewise unique between $n$
and the first Dirichlet zero, since $1+z>0$ at $n$. Each proposed Robin
bracket lies strictly inside its corresponding certified gap. The
location relative to both window endpoints is resolved by strict
inequalities. For all remaining angular orders $n>\rho+\omega$,
positivity up to $n$ excludes any root in the window.

For the second argument put $\mathfrak g_n=J_n+xJ_n'$ and
$\mathfrak d_n=x^2+1-n^2$. On $x>\max(n,1)$, direct differentiation gives
\begin{equation}\label{br:eq:Robin-ode}
 \left(\mathfrak g_n\sqrt{x/\mathfrak d_n}\right)''+
 \left(1+\frac{1/4-n^2}{x^2}+\frac4{\mathfrak d_n}
              -\frac{3x^2}{\mathfrak d_n^2}\right)
 \mathfrak g_n\sqrt{x/\mathfrak d_n}=0.
\end{equation}
The coefficient is at most $21/4$: $\mathfrak d_n>1$, the second term
is at most $1/4$, and the last term is nonpositive. The same spacing
argument gives separation greater than $\pi/\sqrt{21/4}>1$.
Roots are simple because at a root
$\mathfrak g_n'=-(\mathfrak d_n/x)J_n\ne0$; simultaneous vanishing of $\mathfrak g_n$ and
$J_n$ would also force $J_n'=0$. Below $n$ the sign of $\mathfrak g_n$ is positive.
The nonzero endpoint signs at
$x_k=\rho-\omega+k/4$, $0\le k\le8\omega$, therefore determine the number
of roots in every quarter-step cell: one for opposite signs, zero for
equal signs. In a cell crossing $n$ the portion below $n$ contains no
root and the portion above it contains at most one. Each bracket from
the first argument belongs to a distinct sign-changing cell, and all
such cells are represented. This establishes a bijection of the state
lists, including the empty next angular sector.
\end{proof}

\subsection{Matrices and integration}\label{br:sec:matrices}
Write $\mathcal V=q-\rho^2$ and form the exact matrices on $\operatorname{Ran}\PP$
\begin{equation}\label{br:eq:matrices}
 \begin{gathered}
 \mathsf B=\PP\mathcal V\PP,\qquad
 \mathsf B_2=\PP\mathcal V^2\PP,\\
 \mathsf D_{ij}=\int_{\partial\D}d\,\tau_i\chi_i \tau_j\chi_j,
 \qquad \mathsf D_{2,ij}=\int_{\partial\D}d^2\tau_i\chi_i \tau_j\chi_j,\\
 \mathsf H=\diag(\alpha_i^2-\rho^2)-\mathsf B+\mathsf D.
 \end{gathered}
\end{equation}
The plus sign in $\mathsf D$ follows from~\eqref{br:eq:Robin-form}.
The two Gram identities are
\begin{equation}\label{br:eq:grams}
 \mathsf G:=\mathsf B_2-\mathsf B^2
 = (\QQ\mathcal V\PP)^*(\QQ\mathcal V\PP),\qquad
 x^T\mathsf D_2x=\nrm{d\operatorname{Tr}(\PP x)}_2^2.
\end{equation}
Both include all complementary modes; $\mathsf D_2$ is not a truncated
square of $\mathsf D$.

Writing $p_0(w)=\sum_jp_j^{(0)}w^{jm}$, the complex Fourier coefficients
of $q$ are
\[
 q_h(r)=\sum_kp_{k+h}^{(0)}p_k^{(0)}r^{(2k+h)m}\qquad(h\ge0).
\]
Subtract $\rho^2$ from $q_0$ to obtain those of $\mathcal V$, and take
the full two-sided convolution for $\mathcal V^2$. For complex coefficients
$v_h$, the angular factors in the cosine basis are $v_0$ in position
$(0,0)$, $\sqrt2v_j$ in $(0,j)$, and $v_{|i-j|}+v_{i+j}$ for $i,j>0$.
The same rule applies to $d$ and $d^2$.

The radial integrations use $N_G=\brGaussNodes$ Gauss--Legendre nodes.
Disjoint sign brackets of the Legendre polynomial in $(-1,1)$, equal in
number to its degree, enclose all its roots. The weights
$2/[(1-x_i^2)\mathrm{Leg}_{N_G}'(x_i)^2]$ include root uncertainty using
\[
 \max_{[-1,1]}|\mathrm{Leg}_{N_G}''|
 = (N_G-1)N_G(N_G+1)(N_G+2)/8.
\]
Real Bessel arguments are evaluated at exact midpoint arguments and
enlarged by their argument radii, since the integer-order integral
representation gives $|J_n'|\le1$ on the real axis. An exported dyadic
basis value carries its interval evaluation error into every matrix entry.

For the analytic quadrature error use the Bernstein ellipse of parameter
two in $x=2r-1$. It has $|r|\le9/8$ and $|\im r|\le3/8$. Set
\begin{equation}\label{br:eq:Gauss-majorant}
 \begin{split}
 P_e&=\sum_j|p_j^{(0)}|(9/8)^{jm},\qquad V_e=P_e^2+\rho^2,\\
 M_e&=8(1+9/8)(1+V_e^2)e^{3\alpha_{\max}/4}/b_*^2,
 \end{split}
\end{equation}
where $b_*$ is the smallest denominator in~\eqref{br:eq:basis} and
$\alpha_{\max}$ is the largest root in the window. The Bessel integral
gives $|J_n(z)|\le e^{|\im z|}$, so $M_e$ bounds every integrand, including
the angular factors and normalisations. The Chebyshev argument
of~\eqref{r2:eq:Gauss-error} depends only on this majorant and Gauss
exactness; hence the error in each entry is at most
\begin{equation}\label{br:eq:Gauss-error}
 16M_e2^{-2N_G}<\brGaussError.
\end{equation}

For the boundary entries complexify $\phi=m\theta$. On
$|\im\phi|\le\log4$ define
$S_4=\sum_{j>0}4^j|p_j^{(0)}|$ and
$P_{1,4}=\sum_j4^jjm|p_j^{(0)}|$.
The Laurent continuation of $d$ is the half-sum of $wp_0'/p_0$ and its
reflection. For the centre, $S_4<7.762<\rho$; this excludes its poles and gives strip
bound $M_d=P_{1,4}/(\rho-S_4)$. A trapezoidal sum with
$N_d=\brDGrid$ nodes adds to every required Fourier coefficient of $d$
and $d^2$ the error
\begin{equation}\label{br:eq:boundary-matrix-error}
 \frac{4(1+M_d^2)4^{-(N_d-h_{\max})}}{1-4^{-N_d}}.
\end{equation}
Here $h_{\max}$ is the largest Fourier index used by the matrices.
This follows by summing the aliased coefficients, each bounded by its
strip supremum times $4^{-|h|}$.

\subsection{The complement as a quadratic form}\label{br:sec:complement}
Let
\begin{equation}\label{br:eq:complement-constants}
 \gamma_0=\omega(2\rho-\omega),\qquad
 t_Q=\frac{2(1+\sqrt{\gamma_0+\rho^2})}{\gamma_0},\qquad
 v_*=2\rho\sigma_0+\sigma_0^2,\quad d_*=\sigma_1/(\rho-\sigma_0).
\end{equation}
Here $\sigma_0=\sum_{j>0}|p_j^{(0)}|$ and
$\sigma_1=\sum_{j>0}jm|p_j^{(0)}|$. On $\overline\D$,
$\rho-\sigma_0\le|p_0|\le\rho+\sigma_0$ and
$|wp_0'|\le\sigma_1$, so $\nrm{\mathcal V}_\infty\le v_*$ and
$\nrm d_\infty\le d_*$.
On $\operatorname{Ran}\QQ$ set
$D_Q=\QQ(-\Delta_{R,1}-\rho^2)\QQ$. Every omitted state obeys
$|\alpha^2-\rho^2|\ge\gamma_0$: below the window this is
$(\rho-\alpha)(\rho+\alpha)\ge\omega(2\rho-\omega)$, and above the window
the bound is larger. Thus $|D_Q|\ge\gamma_0$.
For $z\in H^1(\D)$ the divergence theorem applied to $x|z|^2$ gives
\begin{equation}\label{br:eq:trace-divergence}
 \nrm{\operatorname{Tr}z}_2^2
 \le2\nrm z_2^2+2\nrm z_2\nrm{\nabla z}_2.
\end{equation}
The estimate extends from smooth functions by $H^1$ approximation.
For $z\in\QQ H^1$, the reference Robin-one form and
$\alpha^2\le|\alpha^2-\rho^2|+\rho^2$ give
\[
 \nrm{\nabla z}_2^2\le
 (1+\rho^2/\gamma_0)\langle|D_Q|z,z\rangle.
\]
Dropping its positive boundary term and inserting this inequality
in~\eqref{br:eq:trace-divergence} proves
\begin{equation}\label{br:eq:trace-energy}
 \nrm{\operatorname{Tr}\QQ|D_Q|^{-1/2}}^2\le t_Q.
\end{equation}
Consequently the complementary form, conjugated by $|D_Q|^{-1/2}$,
is $\operatorname{sgn}D_Q$ plus a bounded self-adjoint perturbation of norm
at most
\begin{equation}\label{br:eq:mu}
 \mu_Q=v_*/\gamma_0+d_*t_Q.
\end{equation}
If $\mu_Q<1$, its inverse in this energy pairing has norm at most
$(1-\mu_Q)^{-1}$. The sign operator is a self-adjoint isometry;
no positivity of $D_Q$ is required.

\begin{proposition}[Robin form Schur bound]\label{br:prop:Schur}
Let $C$ be a real square matrix of dimension $\brStates$ and $\Theta$
a real diagonal matrix with nonzero entries. Suppose
\begin{equation}\label{br:eq:congruence}
 \begin{split}
 \nrm{C^TC-I}_{\rm row}&\le\epsilon_C<1,\\
 \nrm{|\Theta|^{-1/2}(C^T\mathsf HC-\Theta)|\Theta|^{-1/2}}_{\rm row}
 &\le\epsilon.
 \end{split}
\end{equation}
Define
\begin{equation}\label{br:eq:Schur-data}
 \begin{split}
 g_V&=\tr(|\Theta|^{-1}C^T\mathsf GC),\qquad
 g_d=\tr(|\Theta|^{-1}C^T\mathsf D_2C),\\
 \eta_Q&=\frac{(\sqrt{g_V/\gamma_0}+\sqrt{t_Qg_d})^2}{1-\mu_Q},\\
 S_*&=\frac{1+\epsilon_C}{\min_i|\Theta_i|(1-\epsilon-\eta_Q)},\\
 t_P&=\max_j\sum_{i:n_i=jm}\tau_i^2,\qquad
 k_*=v_*/\sqrt{\gamma_0}+d_*\sqrt{t_Qt_P},\\
 \delta_Q&=\frac{k_*}{\sqrt{\gamma_0}(1-\mu_Q)}.
 \end{split}
\end{equation}
If $\mu_Q<1$ and $\epsilon+\eta_Q<1$, the problem~\eqref{br:eq:Robin}
has an $L^2$ inverse with
\begin{equation}\label{br:eq:BLtwo}
 \nrm{L^{-1}}_{2\to2}\le B_{L^2}
 :=(1+\delta_Q^2)S_*+\frac1{\gamma_0(1-\mu_Q)}.
\end{equation}
\end{proposition}
\begin{proof}
The finite spectral projection preserves $H^1$, and the energy space
of $|D_Q|^{1/2}$ is $\QQ H^1$ with an equivalent norm. Thus the
block decomposition of~\eqref{br:eq:Robin-form} is legitimate on
$\operatorname{Ran}\PP\oplus\QQ H^1$. Its finite diagonal block is
$\mathsf H$ and its complementary block is the form just estimated.
Denote the coupling from the finite block into the complementary dual
space by $\mathfrak b$. Its volume part is $-\QQ\mathcal V\PP$; its
boundary part is the pairing with $d\operatorname{Tr}\PP$.
By~\eqref{br:eq:grams} and~\eqref{br:eq:trace-energy}, after multiplication
on the right by $C|\Theta|^{-1/2}$, their energy Hilbert--Schmidt norms
are at most $\sqrt{g_V/\gamma_0}$ and $\sqrt{t_Qg_d}$.
The triangle inequality bounds their sum. Pairing twice with the
complementary inverse yields the relative Schur correction at most
$\eta_Q$.

The matrices in~\eqref{br:eq:congruence} are symmetric, so their maximum
row sums bound their Euclidean operator norms. The first makes $C$
invertible and gives $\nrm C^2\le1+\epsilon_C$. If $\mathsf S$ is the
finite Schur complement, then
$|\Theta|^{-1/2}C^T\mathsf SC|\Theta|^{-1/2}$ is
$\operatorname{sgn}\Theta$ plus a symmetric perturbation of norm at
most $\epsilon+\eta_Q<1$. Its inverse norm is at most
$(1-\epsilon-\eta_Q)^{-1}$, and undoing the congruence gives
$\nrm{\mathsf S^{-1}}\le S_*$.

The finite trace map has squared norm at most $t_P$: angular
orthogonality gives
$\nrm{\operatorname{Tr}\PP x}_2^2
=\sum_j|\sum_{i:n_i=jm}\tau_ix_i|^2\le t_P\nrm x_2^2$.
The unweighted coupling into the complementary energy space therefore
has norm at most $k_*$. The complementary inverse, followed by inclusion
into $L^2$, bounds the off-diagonal solve by $\delta_Q$ and its
$L^2$ diagonal solve by $1/[\gamma_0(1-\mu_Q)]$.
The symmetric block inverse of $-L$ has the form
\[
 \begin{pmatrix}I\\-X_Q\end{pmatrix}\mathsf S^{-1}
 \begin{pmatrix}I&-X_Q^*\end{pmatrix}
 +\begin{pmatrix}0&0\\0&\mathfrak M_Q^{-1}\end{pmatrix},
 \qquad \nrm{X_Q}\le\delta_Q,
\]
where $\mathfrak M_Q^{-1}$ is the complementary form inverse restricted
to $L^2$. This proves~\eqref{br:eq:BLtwo}. The formula constructs a weak
solution for every $L^2$ right side; the invertible complementary and
Schur blocks also prove uniqueness. Smooth coefficients and boundary
give the usual Robin $H^2$ regularity, so this form inverse is the
operator inverse asserted in the proposition.
\end{proof}

For the centre~\eqref{br:eq:input}, the enclosures give
\begin{equation}\label{br:eq:spectral-numbers}
 \begin{gathered}
 d_*<\brDinf<1,\qquad \mu_Q<\brMu,\qquad
 \epsilon_C<\brEpsC,\qquad\epsilon<\brEps,\\
 \min_i|\Theta_i|>\brDiagLo,\qquad
 \eta_Q<\brSchur,\qquad B_{L^2}<\brBLtwo.
 \end{gathered}
\end{equation}
In particular $\beta\ge1-d_*>0$.

\section{Transfer to the working norms}\label{br:sec:transfer}
Define the centre coefficient sums
\begin{equation}\label{br:eq:P}
 \begin{gathered}
 P=\sum_j2^j|p_j^{(0)}|,\qquad
 P_1=\sum_j2^jjm|p_j^{(0)}|,\qquad
 P_2=\sum_j2^j(jm)^2|p_j^{(0)}|,\\
 J_0=(2\rho-P)^{-1},\qquad
 d_0=J_0P_1,\qquad
 d_1=J_0(P_1+P_2)+J_0^2P_1^2.
 \end{gathered}
\end{equation}
The reciprocal is well defined because $P<\brP<2\rho$.
The Neumann series for $1/p_0$ in $W_0$ gives norm at most $J_0$.
For $\mathcal Q_0=wp_0'/p_0$,~\eqref{br:eq:realpart-norm} and differentiation give
\begin{equation}\label{br:eq:d-norms}
 \nrm d_{W_0}\le d_0,\qquad \nrm d_{W_1}\le d_1.
\end{equation}
Indeed $\nrm{wp_0'}_{W_0}=P_1$, while
$\nrm{\partial_\theta \mathcal Q_0}_{W_0}\le J_0P_2+J_0^2P_1^2$.
Also multiplication by $q$ on $Y$ has norm at most $P^2$, and
\begin{equation}\label{br:eq:potential-norm}
 \nrm{q-\rho^2}_Y\le P^2-\rho^2.
\end{equation}
For the latter, remove the constant--constant product in $p_0\overline{p_0}$
and bound $2^{|j-k|}$ by $2^{j+k}$ on every remaining term.

\begin{proposition}[Robin inverse and two trace derivatives]\label{br:prop:transfer}
Let $j_{\rm h}$ be the least positive integer with $j_{\rm h}m>\rho$. Put
\begin{equation}\label{br:eq:high-data}
 \gamma_{\rm h}=(j_{\rm h}m)^2-\rho^2,\qquad
 b_{\rm h}=\frac{P^2-\rho^2}{\gamma_{\rm h}}+\frac{d_0}{1+\sqrt{\gamma_{\rm h}}}.
\end{equation}
Suppose that~\eqref{br:eq:Robin} has an inverse on $L^2$ of the
symmetric sector with norm at most $B_{L^2}$, and that $\beta>0$
and $b_{\rm h}<1$. For $f\in Y$ the Robin solution of
\eqref{br:eq:Robin} belongs to $Y$, its boundary trace belongs to $W_2$,
and
\begin{equation}\label{br:eq:transfer-bounds}
 \nrm V_Y\le B_Y\nrm f_Y,\qquad
 \nrm{\operatorname{Tr}V}_{W_s}\le B_{Ts}\nrm f_Y\quad(s=1,2),
\end{equation}
where, with $P_{\rm un}=\sum_j|p_j^{(0)}|$,
\begin{equation}\label{br:eq:transfer-constants}
 \begin{split}
 \tau_*^2&=2B_{L^2}^2+
 2B_{L^2}\sqrt{P_{\rm un}^2B_{L^2}^2+B_{L^2}},\\
 b_{\rm low}&=\bigl[(1+P^2B_{L^2})/2+\tau_*\bigr]/\sqrt2,\\
 B_Y&=\frac{(1+2\sum_{j=1}^{j_{\rm h}-1}2^j)b_{\rm low}+1/\gamma_{\rm h}}{1-b_{\rm h}},\\
 B_{T1}&=1+P^2B_Y+d_0B_Y,\\
 B_{T2}&=1+P^2B_Y+d_1B_{T1}.
 \end{split}
\end{equation}
\end{proposition}
\begin{proof}
\emph{Energy and low modes.}
The weak equation~\eqref{br:eq:Robin-form}, tested with $V$, gives
\[
 \nrm{\nabla V}_2^2+\int_{\partial\D}\beta V^2
 =\int_\D qV^2-\int_\D fV.
\]
Dropping the positive boundary term and using the $L^2$ inverse yields
\[
 \nrm{\nabla V}_2^2
 \le(P_{\rm un}^2B_{L^2}^2+B_{L^2})\nrm f_2^2.
\]
The divergence estimate~\eqref{br:eq:trace-divergence} then gives
$\nrm{\operatorname{Tr}V}_2\le\tau_*\nrm f_2$.
For each angular coefficient write
$V_n=K_n(f-qV)_n+t_nr^{|n|}$, where $K_n$ is the zero-Dirichlet inverse of the radial operator in
Lemma~\ref{r2:lem:Poisson} and $t_n$ is its boundary trace.
The radial Green kernel in Lemma~\ref{r2:lem:Poisson} is bounded by
$-\log\max(r,s)$ and
$\int_0^1(\log s)^2s\,ds=1/4$. Therefore
\[
 \sup_r|K_n(f-qV)_n|\le\tfrac12\nrm{(f-qV)_n}_{L^2(r\,dr)}.
\]
The three comparisons completing the estimate are
\[
 \nrm f_2\le\nrm f_Y/\sqrt2,\qquad
 \nrm{(f-qV)_n}_{L^2(r\,dr)}\le\nrm{f-qV}_2,
 \qquad |t_n|\le\nrm{\operatorname{Tr}V}_2.
\]
The first follows from Parseval and $\int_0^1r\,dr=1/2$, the second
from orthogonality, and the third from the boundary Fourier coefficient
formula. Thus every coefficient has supremum at most
$b_{\rm low}\nrm f_Y$. The low part $V_{\rm low}$, consisting of $|n|<j_{\rm h}m$,
has weighted norm at most
$(1+2\sum_{j=1}^{j_{\rm h}-1}2^j)b_{\rm low}\nrm f_Y$.

\emph{Construction of the high modes.}
For $|n|\ge j_{\rm h}m$, the operator $-\Delta_n-\rho^2$ has positive radial
potential $n^2/r^2-\rho^2\ge\gamma_{\rm h}$. The maximum principle with
Robin-one boundary data gives interior-forcing norm at most
$1/\gamma_{\rm h}$. This may also be seen by comparing with the constant
supersolution $\nrm{\mathfrak f}_\infty/\gamma_{\rm h}$ for the interior forcing
$\mathfrak f$ of~\eqref{br:eq:high-equations}. The supersolution has
positive Robin boundary datum; apply the comparison to both signs of the real solution.
For a pure boundary source the regular homogeneous solution is a multiple
of $J_{|n|}(\rho r)$, which is positive and increasing because $|n|>\rho$.
Its logarithmic derivative
$z_n=r\rho J_{|n|}'(\rho r)/J_{|n|}(\rho r)$ obeys
\[
 rz_n'=n^2-\rho^2r^2-z_n^2.
\]
The power series at zero gives
$z_n=|n|-\rho^2r^2/[2(|n|+1)]+O(r^4)$, greater than
$\sqrt{n^2-\rho^2r^2}$ for small positive $r$. At a first contact with
this decreasing comparison function, $z_n'=0$ and the difference has
positive derivative, excluding the contact. Consequently the solution
with boundary source $\varsigma_n$ has supremum at most
$|\varsigma_n|/(1+\sqrt{\gamma_{\rm h}})$.

Let $\Pi_H$ be the angular high projection. In the high subspace of $Y$
solve by contraction
\begin{equation}\label{br:eq:high-equations}
 \begin{split}
 (\Delta+\rho^2)\widetilde V_H
 &=\Pi_H\bigl(f-(q-\rho^2)(V_{\rm low}+\widetilde V_H)\bigr),\\
 (\widetilde V_H)_r+\widetilde V_H
 &=-\Pi_H\bigl(d\operatorname{Tr}(V_{\rm low}+\widetilde V_H)\bigr).
 \end{split}
\end{equation}
The interior and boundary estimates, the trace norm $Y\to W_0$ at most
one, and~\eqref{br:eq:d-norms} give Lipschitz constant $b_{\rm h}<1$.
This constructs $\widetilde V_H$ with finite $Y$ norm before any total
analytic norm is assumed.

For its identification with $\Pi_HV$ we first justify the energy
equation. In each mode the constructed solution is regular at zero and
solves its radial ODE with Robin-one boundary source. If $\mathfrak f_n$ and $\varsigma_n$
are the right sides of~\eqref{br:eq:high-equations}, the mode energy
identity is
\[
 \int_0^1\bigl(|(\widetilde V_{H,n})'|^2+
 (n^2/r^2-\rho^2)|\widetilde V_{H,n}|^2\bigr)r\,dr
 +|\widetilde V_{H,n}(1)|^2
 =-\re\int_0^1\mathfrak f_n\overline{\widetilde V_{H,n}}r\,dr
 +\re(\varsigma_n\overline{\widetilde V_{H,n}(1)}).
\]
The weighted supremum sums for the solution and $\mathfrak f_n$, and the weighted
boundary sum for $\varsigma_n$, make the sum of the absolute right sides finite.
Adding the finite $\rho^2L^2$ term shows convergence of the mode series
in $H^1$. The mode equations, first tested against finite angular sums
and then by $H^1$ density, give the weak high-subspace equation and its
boundary pairing. In particular the boundary source is an actual weak
Robin datum.

The original $L^2$ inverse gives $V\in H^2$; testing its weak equation
against a radial test function times $\cos(n\theta)$ gives
\begin{equation}\label{br:eq:mode-Robin}
 V_n'(1)=-(\beta\operatorname{Tr}V)_n.
\end{equation}
Here $V_n$ is $C^2$ on $(0,1]$ since $(f-qV)_n$ is continuous.
Thus $V_H=\Pi_HV$ satisfies the same weak projected equation as
$\widetilde V_H$, with its own occurrence on the right side. Their
difference $z\in\Pi_HH^1$ satisfies
\[
 \nrm{\nabla z}_2^2+\int_{\partial\D}\beta z^2=\int_\D qz^2.
\]
Angular orthogonality gives
$\nrm{\nabla z}_2^2\ge(j_{\rm h}m)^2\nrm z_2^2$.
The inequality $b_{\rm h}<1$ implies $(j_{\rm h}m)^2>P^2\ge\nrm q_\infty$.
Since $\beta>0$, $z=0$. This identifies the analytic high solution
with the original Robin solution.

The high fixed-point equation now gives
$\nrm{V_H}_Y\le\nrm f_Y/\gamma_{\rm h}+b_{\rm h}\nrm V_Y$.
Add the low bound and rearrange to obtain $B_Y$.

\emph{Two boundary derivatives.}
Integrate the radial equation $\Delta_nV_n=(f-qV)_n$ against
$r^{|n|+1}$. The moment identity underlying
Lemma~\ref{r2:lem:Poisson}, now with nonzero Dirichlet trace, reads
\[
 \int_0^1r^{|n|+1}(f-qV)_n\,dr=V_n'(1)-|n|t_n.
\]
Using~\eqref{br:eq:mode-Robin} gives
\begin{equation}\label{br:eq:trace-identity}
 (|n|+1)t_n=-\int_0^1r^{|n|+1}(f-qV)_n\,dr-(dt)_n.
\end{equation}
Since $t\in W_0$ already, summing this formula with weights $R^{|n|}$
first proves $t\in W_1$. Its norm is at most
$\tfrac12\nrm{f-qV}_Y+d_0\nrm t_{W_0}$, bounded by the stated $B_{T1}$.
Multiply the same identity by $|n|+1$ and use
$(|n|+1)/(|n|+2)\le1$. The $W_1$ algebra bound on $dt$ then proves
$t\in W_2$ and the stated $B_{T2}$.
\end{proof}

The centre enclosures give $b_{\rm h}<\brHighFactor<1$ and
\begin{equation}\label{br:eq:transfer-numbers}
 B_Y<\brBY,\qquad B_{T1}<\brBTone,\qquad B_{T2}<\brBTtwo.
\end{equation}
The second boundary derivative is needed because reconstruction of $p$
from its primitive consumes one derivative while the unknown shape
belongs to $W_1$.
\section{All-mode residual and the exact linearisation inverse}\label{br:sec:inverse}
\subsection{Strip bounds and the clamped residual}\label{br:sec:residual}
Complexify $\phi=m\theta$ to $|\im\phi|\le\log4$. Set
\begin{equation}\label{br:eq:strip-data}
 \begin{gathered}
 S_c=\sum_{j>0}4^j|c_j^{(0)}|,\qquad
 \delta_s=(2\rho S_c+S_c^2)/\rho,\qquad
 \varepsilon_s=\tfrac12\log\frac{\rho+S_c}{\rho-S_c},\\
 P_4=\sum_j4^j(jm+1)|c_j^{(0)}|.
 \end{gathered}
\end{equation}
For $0\le r\le1$ put
$A_\pm=\rho+\sum_{j>0}c_j^{(0)}r^{jm}e^{\pm ij\phi}$.
For the centre, $S_c<0.0571$, so $\rho>2S_c$; this places each $A_\pm$ in the disc
$|A-\rho|\le S_c<\rho/2$. Define
\begin{equation}\label{br:eq:strip-root}
 B_s=\sqrt{A_+}\sqrt{A_-}
\end{equation}
using the principal roots of the individual factors. Both roots are
analytic, and $|\arg B_s|<\pi/6$, hence $\re B_s>0$ and
$|B_s+\rho|\ge\rho$. Therefore
\begin{equation}\label{br:eq:strip-geometry}
 |B_s-\rho|=\frac{|A_+A_--\rho^2|}{|B_s+\rho|}\le\delta_s,
 \qquad |A_\pm/B_s|\le e^{\varepsilon_s}.
\end{equation}

The complexified physical radius is $rB_s$. The positive angular branch
of order $n=jm$ is
$J_n(rB_s)e^{ij\phi}(A_+/B_s)^n$, with the reflected expression for
the negative branch. The entire Cartesian Bessel power series interprets
this formula also at $r=0$. For integer order the contour shift
in~\eqref{r2:eq:Bessel-contour} uses only the Bessel integral and gives
\[
 |J_n(z)|\le
 \exp(-n\vartheta+|\re z|\sinh\vartheta+|\im z|\cosh\vartheta)
 \quad(\vartheta\ge0).
\]
Choose $\vartheta_j=0$ for $jm\le\rho+\delta_s$ and otherwise
$\vartheta_j=\acosh(jm/(\rho+\delta_s))$. Define
\begin{equation}\label{br:eq:strip-summands}
 \mathfrak u_j=\frac{|a_j|}{s_j}4^j
 \exp\bigl(jm\varepsilon_s-jm\vartheta_j
       +(\rho+\delta_s)\sinh\vartheta_j+\delta_s\cosh\vartheta_j\bigr).
\end{equation}
Since $|\re(rB_s)|\le\rho+\delta_s$ and $|\im(rB_s)|\le\delta_s$,
$\sum_j\mathfrak u_j$ bounds $U_*$ uniformly on the strip. The coefficient
bound $M4^{-|j|}$ yields the weighted sum factor
$1+2\sum_{j\ge1}2^{-j}=3$. Consequently
\begin{equation}\label{br:eq:field-bounds}
 \begin{split}
 \nrm{U_*}_Y&\le U_b:=3\sum_j\mathfrak u_j,\\
 \nrm{\partial_wU_*}_{\Yc},\ \nrm{\partial_{\bar w}U_*}_{\Yc}
 &\le \mathfrak D_U:=3R4^{2/m}P_4e^{\varepsilon_s}
                      \sum_j\mathfrak u_je^{\vartheta_j}.
 \end{split}
\end{equation}
For the derivative bound use the Wirtinger recurrences displayed in
Subsection~\ref{r2:sec:centre}: a physical derivative changes Bessel
order by one, costs at most $e^{\vartheta_j}$ in the contour bound, and
changes the angular factor by at most $4^{1/m}e^{\varepsilon_s}$.
The chain rule contributes $P_4$. The resulting angular frequencies are
$km\pm1$; multiplying by $e^{\mp i\theta}$ makes them periodic in
$\phi$, and $R^{|km\pm1|}\le R2^{|k|}$ costs the further factors
$R4^{1/m}$. These estimates give the enlarged $\mathfrak D_U$ in
\eqref{br:eq:field-bounds}, including the order-zero term via
the Bessel identity $J_{-1}(z)=-J_1(z)$.

The common strip majorant for both defects in~\eqref{br:eq:defects} is
\begin{equation}\label{br:eq:trace-majorant}
 M_{\rm tr}=U_b+2\mathfrak D_U+P_4.
\end{equation}
Indeed the radial derivative is
$e^{i\theta}U_w+e^{-i\theta}U_{\bar w}$, whose strip bound is at most
$2\mathfrak D_U$ by the estimates just derived. The complexification of $|p_0|$
is the product of the principal square roots of $p_0(e^{i\theta})$ and
its reflected series. The inequality $S_4<\rho$ of Subsection~\ref{br:sec:matrices} makes this product analytic on the strip, with modulus at most
$P_4$. This accounts for the additional term in~\eqref{br:eq:trace-majorant}.

Use $N_b=\brTraceGrid$ interval trapezoidal nodes and retain the cosine
coefficients through $L_b=\brTraceCutoff<N_b/2$. Summing the aliases gives
the common retained-coefficient error
\begin{equation}\label{br:eq:trace-alias}
 e_b=\frac{8M_{\rm tr}4^{-(N_b-L_b)}}{1-4^{-N_b}}.
\end{equation}
For $j>L_b$, $|D_j|,|N_j|\le2M_{\rm tr}4^{-j}$.
Cauchy's estimate gives the same bound for every $j\ge1$, and the
cosine coefficients of $|p_0|$ are bounded by $2P_44^{-j}$. Thus the
series~\eqref{br:eq:centre} for $\ell$ and the series for $\Bc(|p_0|)$
converge with all derivatives on every closed disc of radius less than
$4^{1/m}$. The complete moments are
\begin{equation}\label{br:eq:trace-moments}
 D^{[k]}=\sum_{j\ge0}2^j(jm+1)^k|D_j|,\qquad
 N^{[k]}=\sum_{j\ge0}2^j(jm+1)^k|N_j|\quad(0\le k\le3).
\end{equation}
Complex Fourier coefficients and one-sided cosine coefficients give the
same norm here, since the two half-coefficients at opposite frequencies
sum to the full cosine coefficient. With $J=L_b+1$, the tail added to
each moment is at most
\begin{equation}\label{br:eq:trace-tail}
 2M_{\rm tr}2^{-J}\sum_{i=0}^k\binom{k}{i}
 (mJ+1)^{k-i}m^i\mathfrak m_i,
 \qquad (\mathfrak m_0,\mathfrak m_1,\mathfrak m_2,\mathfrak m_3)
 =(2,2,6,26).
\end{equation}
These four integers are the exact sums
$\sum_{t\ge0}t^i2^{-t}$, as in~\eqref{r2:eq:tail-sums}. Expanding
$(m(J+t)+1)^k$ proves~\eqref{br:eq:trace-tail} and includes every
uncomputed harmonic.

The lift polynomials are identical to those of the Schiffer centre;
their norm estimates depend only on the defect coefficients. In the
present notation they give
\begin{equation}\label{br:eq:lift-constants}
 \begin{aligned}
 L_0&=D^{[0]}+N^{[0]}/2,&
 L_1&=2R(D^{[2]}+N^{[1]}),\\
 L_\Delta&=2(D^{[2]}+N^{[1]}),&
 L_{\Delta,1}&=2R(D^{[3]}+N^{[2]}).
 \end{aligned}
\end{equation}
They bound $\ell$, each Wirtinger derivative of $\ell$, $\Delta\ell$,
and each Wirtinger derivative of $\Delta\ell$, respectively.
For example,
\[
 \begin{split}
 \Delta(h_n^D\cos n\theta)&=-2n(n+1)r^n\cos n\theta,\\
 \Delta(h_n^N\cos n\theta)&=2(n+1)r^n\cos n\theta.
 \end{split}
\]
Writing the lifts in $w,\bar w$ proves the
derivative bounds with one factor $R$ for the shifted angular frequency;
the displayed enlargements include $n=0$.
Since $\nrm{\partial_wq}_{\Yc}\le RPP_1$ and multiplication by $q$
has norm at most $P^2$, equation~\eqref{br:eq:E} gives
\begin{equation}\label{br:eq:residual-bounds}
 \begin{split}
 \nrm E_Y&\le E_0:=L_\Delta+P^2L_0,\\
 \nrm{\partial_wE}_{\Yc},\ \nrm{\partial_{\bar w}E}_{\Yc}
 &\le E_1:=L_{\Delta,1}+RPP_1L_0+P^2L_1,\\
 \nrm{\operatorname{Tr}E}_{W_1}&\le E_{W1}:=E_0+2RE_1.
 \end{split}
\end{equation}
The last bound uses $\partial_\theta E=i(wE_w-\bar wE_{\bar w})$.
Likewise $\Delta U_0=-qU_*-\Delta\ell$ implies
\begin{equation}\label{br:eq:Lambda}
 \nrm{\Delta U_0}_Y\le\Lambda_0:=P^2U_b+L_\Delta,
 \qquad
 \nrm{\partial_w\Delta U_0}_{\Yc}\le
 \Lambda_1:=RPP_1U_b+P^2\mathfrak D_U+L_{\Delta,1},
\end{equation}
with the same bound for the conjugate derivative. In particular,
\begin{equation}\label{br:eq:residual-numbers}
 U_b<\brUb,\qquad \mathfrak D_U<\brDUb,\qquad
 E_0<\brEzero,\qquad E_1<\brEone.
\end{equation}

\subsection{Material derivative and boundary identities}\label{br:sec:linearization}
All quantities below are evaluated at $x_0=(g_0,p_0)$. For a variation
$(\delta g,k)\in X$ define
\begin{equation}\label{br:eq:variation}
 \eta(w)=\int_0^wk(z)\,dz,\qquad
 v=\eta/p_0=wa,\qquad
 T=v\partial_wU_0+\bar v\partial_{\bar w}U_0.
\end{equation}
The symbol $a$ denotes only the holomorphic real-coefficient shape
function in this formula; the $a_j$ in~\eqref{br:eq:input} are the fixed
field coefficients. Put $b=|p_0|$ on the circle and denote its real
derivative by $\mathrm Db[k]=b\re(k/p_0)$.
Differentiation of~\eqref{br:eq:F} yields
\begin{equation}\label{br:eq:DF}
 DF(x_0)(\delta g,k)=LH+\delta q\,U_0,
 \qquad H=K\delta g+\Bc(\mathrm Db[k]),\quad
 \delta q=2\re(\bar p_0k).
\end{equation}
The material identity of the Schiffer part uses only holomorphicity
and the equation $\Delta U_0+qU_0=E$, so it gives here
\begin{equation}\label{br:eq:material}
 LT+\delta q\,U_0=\Ecal(k),\qquad
 \Ecal(k)=vE_w+\bar vE_{\bar w}+(v'+\bar v')E.
\end{equation}
To see all cancellations, $\Delta(vU_w)=v\partial_w\Delta U+v'\Delta U$,
and the conjugate formula gives
\begin{equation}\label{br:eq:DeltaT}
 \Delta T=v\partial_w\Delta U_0+\bar v\partial_{\bar w}\Delta U_0
                  +(v'+\bar v')\Delta U_0.
\end{equation}
The identity $k=(p_0v)'$ gives
$\delta q=vq_w+\bar vq_{\bar w}+(v'+\bar v')q$.
Substitute $\Delta U_0=E-qU_0$ in~\eqref{br:eq:DeltaT}; the terms involving
derivatives of $q$ cancel with $\delta qU_0$, and the other terms involving
$U_0$ cancel with $qT$. This proves~\eqref{br:eq:material}.

The boundary Hessian calculation differs from the Schiffer case.
Exact clamping gives $U_0=0$, $(U_0)_r=b$, $(U_0)_{r\theta}=b_\theta$,
and $(U_0)_{rr}=E-b$ on the circle. Since
$T=r\re a\,(U_0)_r+\im a\,(U_0)_\theta$, radial differentiation gives
\begin{equation}\label{br:eq:T-boundary}
 T=b\re a,\qquad
 T_r=b\re(a+wa')+(E-b)\re a+b_\theta\im a.
\end{equation}
For $\mathcal Q_0=wp_0'/p_0$ one has
\[
 \mathrm Db[k]=b\re\bigl((1+\mathcal Q_0)a+wa'\bigr),\qquad
 b_\theta=-b\im \mathcal Q_0.
\]
The tangential terms cancel in the difference, leaving
\begin{equation}\label{br:eq:boundary-difference}
 \mathrm Db[k]-T_r=(\beta-E/b)T,
 \qquad \beta=1+\re \mathcal Q_0.
\end{equation}
Thus the residual introduces a boundary term as well as the interior
term $\Ecal$. Define its Neumann lift and the corrected linearisation by
\begin{equation}\label{br:eq:corrected}
 \mathcal C(k)=\Bc(E\re a),\qquad
 \widetilde D=DF(x_0)-\Ecal+L\mathcal C.
\end{equation}
Equations~\eqref{br:eq:DF} and~\eqref{br:eq:material} then give
$\widetilde D(\delta g,k)=L(H-T+\mathcal C(k))$.

\subsection{The exact two-sided inverse}\label{br:sec:two-sided}
\begin{proposition}\label{br:prop:inverse}
Suppose that~\eqref{br:eq:Robin} has the inverse and trace bounds of
Proposition~\ref{br:prop:transfer}, and that $P<2\rho$. Then
$\widetilde D\colon X\to Y$ of~\eqref{br:eq:corrected} has a bounded
two-sided inverse $A\colon Y\to X$ with $\nrm A\le A_*$, where $A_*$
is defined in~\eqref{br:eq:inverse-norm-data}.
\end{proposition}
\begin{proof}
\emph{Right inverse.}
Given $f\in Y$, solve~\eqref{br:eq:Robin} by the Robin inverse and let
$t=\operatorname{Tr}V\in W_2$. Since $b$ does not vanish, prescribe
\begin{equation}\label{br:eq:shape-recovery}
 \re a=-t/b.
\end{equation}
The reciprocal estimates in Subsection~\ref{br:sec:inverse-norm} below
give $1/b\in W_2$. The right side is real, even and has frequencies in
$m\Z$. If its cosine coefficients are $\widehat a_j$, define
$a(w)=\sum_{j\ge0}\widehat a_jw^{jm}$. Thus the cosine coefficients of
$-t/b$ are the Taylor coefficients of $a$. Equivalently retain the
constant complex Fourier coefficient and double every positive one.
This is the unique solution in the holomorphic real-coefficient class;
the imaginary constant of the general real-part problem is excluded.
Equation~\eqref{br:eq:realpart-norm} gives equality of the $W_2$ norms.
Define, in this order,
\begin{equation}\label{br:eq:inverse-recipe}
 v=wa,\quad k=(p_0wa)',\quad
 H=V+T-\mathcal C(k),\quad
 \delta g=\Delta\bigl(H-\Bc(\mathrm Db[k])\bigr).
\end{equation}
The boundary value of $H$ is $t+b\re a=0$. Using
$V_r=-\beta t=\beta T$ and~\eqref{br:eq:boundary-difference},
\[
 H_r=\beta T+T_r-E\re a=\mathrm Db[k].
\]
These equalities hold mode by mode by~\eqref{br:eq:mode-Robin}.
The function $H-\Bc(\mathrm Db[k])$ has zero Cauchy data and a Laplacian
in $Y$. The radial moment identity therefore gives $\delta g\in X_g$,
and Poisson uniqueness gives
$K\delta g=H-\Bc(\mathrm Db[k])$.
The interior identity now gives
$\widetilde D(\delta g,k)=LV=f$. The recipe is linear in $f$.

\emph{Derivative count and regularity.}
For an arbitrary $k\in W_1$,
\begin{equation}\label{br:eq:derivative-count}
 \nrm{\eta/w}_{W_2}=\nrm k_{W_1},\qquad
 \nrm{(p_0wa)'}_{W_1}=\nrm{p_0a}_{W_2}.
\end{equation}
Both identities follow coefficient by coefficient from multiplication
or division by $jm+1$. The Robin solve gains two boundary derivatives,
division by $b$ preserves $W_2$, and reconstruction of $k$ spends one.
The Laplacian of the Neumann lift spends one boundary derivative,
which is available from $k\in W_1$. All holomorphic shape functions
are smooth near $\overline\D$ because $R>1$. The Robin solution is a
weak $H^2$ solution, with continuous right side; standard elliptic
regularity gives $W^{2,s}$ for finite $s$ and $C^1$ boundary traces.
The mode-wise boundary identity proved above justifies the zero
Neumann trace of $H-\Bc(\mathrm Db[k])$. The estimates below show that
its Laplacian belongs to $Y$, and hence that the recipe maps $Y$ to $X$.

\emph{Left inverse.}
Start with any $(\delta g,k)\in X$. Form $\eta,v,a,T$ by
\eqref{br:eq:variation}, $H$ by~\eqref{br:eq:DF}, and put
$V=H-T+\mathcal C(k)$. The apparent singularity in
$a=\eta/(wp_0)$ at zero is removable, and
\eqref{br:eq:derivative-count} gives $a\in W_2$.
Its Dirichlet trace is $-b\re a$, and
\eqref{br:eq:boundary-difference} gives $V_r+\beta V=0$.
Also $LV=\widetilde D(\delta g,k)$ distributionally.
The terms involving the fixed smooth centre and the holomorphic shape
are smooth near the circle; $K\delta g$ has $C^1$ traces. Thus this
function is the weak Robin solution. Uniqueness recovers the same $V$
when the right-inverse recipe is applied to $\widetilde D(\delta g,k)$.
Its Dirichlet trace divided by $-b$ recovers $\re a$, and the unique
real-coefficient real-part solution recovers $a$. Then $k$, $T$,
$\mathcal C(k)$, $H$ and $\delta g$ are recovered in turn.
Consequently $A\widetilde D=I_X$ and $\widetilde DA=I_Y$.
\end{proof}

\subsection{Inverse norm and left Newton defect}\label{br:sec:inverse-norm}
The reciprocal bounds required above are
\begin{equation}\label{br:eq:reciprocals}
 J_1=J_0+P_1J_0^2,\qquad
 J_2=J_0+(2P_1+P_2)J_0^2+2P_1^2J_0^3.
\end{equation}
They bound $1/p_0$ in $W_1$ and $W_2$, respectively; $J_2$ also bounds
$1/b$ in $W_2$. To prove the holomorphic bounds differentiate the
reciprocal once and twice in the angular variable and use
$\nrm h_{W_2}=\nrm h_{W_0}+2\nrm{\partial_\theta h}_{W_0}
+\nrm{\partial_\theta^2h}_{W_0}$. For the boundary reciprocal,
$\log b=(\log p_0+\log\overline{p_0})/2$ has first derivative norm at
most $J_0P_1$ and second derivative norm at most
$J_0P_2+J_0^2P_1^2$. The product of the individual square-root reciprocal
series gives $\nrm{1/b}_{W_0}\le J_0$; differentiating
$\exp(-\log b)$ gives exactly the stated $J_2$ enlargement.

With $\delta_p=(P-\rho)/\rho<1$, set
\begin{equation}\label{br:eq:b-bound}
 B_{01}=\rho(1+\delta_p)^2+
                 \frac{1+\delta_p}{\sqrt{1-\delta_p}}P_1.
\end{equation}
This bounds $\nrm b_{W_1}$. Indeed the sum of absolute nonconstant
coefficients in the series for $\sqrt{1+z}$ at $|z|=\delta_p$ is
$1-\sqrt{1-\delta_p}\le\delta_p$. Apply this to $\sqrt{p_0/\rho}$ and
its reflection and multiply. Differentiating a square root adds at
most $P_1/[2\sqrt\rho\sqrt{1-\delta_p}]$; the two product derivatives
give the second term of~\eqref{br:eq:b-bound}. In particular,
\begin{equation}\label{br:eq:Db-bound}
 \nrm{\mathrm Db[k]}_{W_1}\le B_{01}J_1\nrm k_{W_1},
\end{equation}
by~\eqref{br:eq:realpart-norm} applied to $k/p_0$.

Define
\begin{equation}\label{br:eq:inverse-norm-data}
 \begin{split}
 B_a&=J_2B_{T2},\qquad B_p=(P+2P_1+P_2)B_a,\\
 B_g&=1+P^2B_Y+2(R\Lambda_1+\Lambda_0+E_{W1})B_a
                         +2B_{01}J_1B_p,\\
 A_*&=B_g+B_p.
 \end{split}
\end{equation}
For $Af=(\delta g,k)$, equation~\eqref{br:eq:shape-recovery} bounds
$\nrm a_{W_2}\le B_a\nrm f_Y$, and the exact second identity
of~\eqref{br:eq:derivative-count} gives
$\nrm k_{W_1}\le B_p\nrm f_Y$ because
$\nrm{p_0}_{W_2}=P+2P_1+P_2$.
As in~\eqref{r2:eq:v-estimates}, $\nrm v_{\Yc}\le R\nrm a_{W_1}$
and $\nrm{v'}_{\Yc}=\nrm a_{W_1}$.
The product identity~\eqref{br:eq:DeltaT} therefore gives
\[
 \nrm{\Delta T}_Y\le2(R\Lambda_1+\Lambda_0)\nrm a_{W_1}.
\]
Also $\nrm{\Delta\mathcal C(k)}_Y\le2E_{W1}\nrm a_{W_1}$.
Applying these estimates and~\eqref{br:eq:Db-bound} to
\eqref{br:eq:inverse-recipe}, with $\Delta V=f-qV$, proves
$\nrm{\delta g}_Y\le B_g\nrm f_Y$ and $\nrm A\le A_*$.

For an arbitrary input shape,
$\nrm a_{W_2}\le J_2\nrm k_{W_1}$ by the first identity of
\eqref{br:eq:derivative-count}. Hence
\begin{equation}\label{br:eq:Newton-data}
 \begin{split}
 \nrm{AF(x_0)}_X&\le Y_0:=A_*E_0,\\
 \nrm{I_X-A DF(x_0)}_{X\to X}&\le Z,\\
 Z&:=A_*J_2\bigl[2(RE_1+E_0)+2E_{W1}+P^2E_0/2\bigr].
 \end{split}
\end{equation}
Indeed $I_X-A DF=A(-\Ecal+L\mathcal C)$, the material correction has
norm at most $2(RE_1+E_0)\nrm a_{W_1}$, and
$\nrm{L\mathcal C}_Y\le(2E_{W1}+P^2E_0/2)\nrm a_{W_1}$.
The last term includes $q\mathcal C$ as well as its Laplacian.
Both identities of Proposition~\ref{br:prop:inverse} are used: the left
identity gives this defect on $X$, and the right identity makes $A$
injective, so a fixed point of the Newton map is a zero of $F$.

\subsection{Symmetry directions and the fixed parameter}\label{br:sec:kernels}
At a true solution, the Eulerian derivative $\dot u$ associated with
normal displacement $h_\nu$ has boundary data
$\dot u=-h_\nu$ and $\partial_\nu\dot u=\kappa h_\nu$.
Thus $\partial_\nu\dot u+\kappa\dot u=0$, agreeing with
\eqref{br:eq:Robin}. Translations have angular character $\pm1$ and
are absent from the $D_m$ invariant sector. The rotation derivative is
reflection-odd and is absent from its real even subspace.
The real conformal coefficients and $\psi(0)=0$, $\psi'(0)>0$ fix the
conformal normalisation.
The constant coefficient $\widehat a_0$ of $a$ generates the radial
reparametrisation $w\mapsto(1+\varepsilon\widehat a_0)w$ of the parameter
disc. It is retained in the inverse and equals the mean of $-t/b$
in~\eqref{br:eq:shape-recovery}. Euclidean dilation is not a homogeneous kernel with
eigenvalue and flux fixed. The flux-preserving family
$u_s(x)=s\,u(x/s)$ has parameter $s^{-2}$; its derivative
$u-x\cdot\nabla u$ satisfies $(\Delta+1)(u-x\cdot\nabla u)=2u$.
Multiplying the field by a constant changes its prescribed flux.
The absence of a kernel within the working sector follows from the
Robin inverse and the two-sided construction above.
\section{Existence, geometry and earlier results}\label{br:sec:existence}
\subsection{The full nonlinear bound}\label{br:sec:nonlinear}
Fix $r_*>0$ with $J_1r_*<1/4$ and consider the closed $X$ ball about
$x_0$ of that radius. Put
\begin{equation}\label{br:eq:ball-data}
 \begin{gathered}
 P_b=P+r_*,\qquad B_b=4B_{01},\qquad J_b=2J_1,\qquad U_c=U_b+L_0,\\
 U_{\rm ball}=U_c+r_*/4+B_bJ_br_*/2.
 \end{gathered}
\end{equation}
On the ball, $\nrm p_{W_0}\le P_b$, $\nrm{|p|}_{W_1}\le B_b$,
$\nrm{1/p}_{W_1}\le J_b$, and $\nrm{U(g,p)}_Y\le U_{\rm ball}$.
For the reciprocal and square root write $p=p_0(1+k/p_0)$.
The $W_1$ algebra gives $\nrm{k/p_0}_{W_1}<1/4$; its Neumann series
bounds the reciprocal by $(4/3)J_1\le J_b$. Multiplying the square-root
series of $1+k/p_0$ and its reflection bounds the modulus by
$4B_{01}=B_b$. Along the segment from $p_0$ to $p$, the formula
$\mathrm D|p|[k]=|p|\re(k/p)$ bounds its change in $W_1$ by
$B_bJ_b\nrm k_{W_1}$. The lift bound and $\nrm K\le1/4$ give the last
bound in~\eqref{br:eq:ball-data}.

\begin{lemma}\label{br:lem:nonlinear}
Throughout the ball, $F$ has second derivative bound
\begin{equation}\label{br:eq:Mtwo}
 \nrm{F''}_{X\times X\to Y}\le M_2
 :=(4+P_b^2)B_bJ_b^2+2P_bB_bJ_b+P_b+2U_{\rm ball}.
\end{equation}
\end{lemma}
\begin{proof}
Differentiating the modulus derivative once more and applying the
$W_1$ product inequality gives
$\nrm{\mathrm D^2|p|[k_1,k_2]}_{W_1}
\le2B_bJ_b^2\nrm{k_1}_{W_1}\nrm{k_2}_{W_1}$.
Let $q_p=|p|^2$ for the varying shape; $q$ continues to denote the fixed
centre potential $|p_0|^2$. For $x_i=(g_i,k_i)$ the full second derivative is
\begin{equation}\label{br:eq:Fsecond}
 \begin{split}
 F''[x_1,x_2]={}&(\Delta+|p|^2)\Bc(\mathrm D^2|p|[k_1,k_2])\\
 &+\mathrm Dq_p[k_1]\bigl(\Bc(\mathrm D|p|[k_2])+Kg_2\bigr)\\
 &+\mathrm Dq_p[k_2]\bigl(\Bc(\mathrm D|p|[k_1])+Kg_1\bigr)
   +\mathrm D^2q_p[k_1,k_2]U.
 \end{split}
\end{equation}
Here $\nrm{\mathrm Dq_p[k]}_Y\le2P_b\nrm k_{W_1}$ and
$\nrm{\mathrm D^2q_p[k_1,k_2]}_Y\le2\nrm{k_1}_{W_1}\nrm{k_2}_{W_1}$.
The first line of~\eqref{br:eq:Fsecond} is bounded by
$(4+P_b^2)B_bJ_b^2$, using both bounds
in~\eqref{br:eq:neumann-lift-bounds}. The two shape/lift products give
$2P_bB_bJ_b$, and the two shape/Poisson products are bounded by $P_b$
using $\nrm K\le1/4$. The final product gives $2U_{\rm ball}$.
For unit $X$ directions these enlarged bounds sum to~\eqref{br:eq:Mtwo}.
\end{proof}

\subsection{The radii polynomial}\label{br:sec:contraction}
Set $C_2=A_*M_2/2$. The Newton map $\mathcal T(x)=x-AF(x)$ satisfies,
for $\nrm{x-x_0}_X\le r_*$,
\begin{equation}\label{br:eq:Newton-map-bounds}
 \nrm{\mathcal T(x)-x_0}_X\le Y_0+Zr_*+C_2r_*^2,
 \qquad \nrm{D\mathcal T(x)}\le Z+2C_2r_*.
\end{equation}
These follow from Taylor's integral remainder,~\eqref{br:eq:Newton-data}
and Lemma~\ref{br:lem:nonlinear}. Thus
\begin{equation}\label{br:eq:radii}
 Y_0+(Z-1)r_*+C_2r_*^2<0,
 \qquad Z+2C_2r_*<1
\end{equation}
make the ball invariant and the map a contraction. The fixed-point
argument in Lemma~\ref{r2:lem:contraction} applies with the quadratic
Taylor majorant~\eqref{br:eq:Newton-map-bounds} in place of the cubic
expansion. The contraction has a unique fixed point;
there $AF(x)=0$. The identity $\widetilde DA=I_Y$ implies $F(x)=0$.
Every zero in the ball is a fixed point, so uniqueness is uniqueness of
the zero of~\eqref{br:eq:F}.

Take $r_*=\brRadius$. The bounds in Table~\ref{br:tab:constants} give
both strict inequalities in~\eqref{br:eq:radii}.
\begin{table}[htbp]
\centering
\begin{tabular}{ll}
\toprule
Quantity & Enclosed bound\\
\midrule
$B_{L^2}$ & $<\brBLtwo$\\
$B_Y$ & $<\brBY$\\
$B_{T1}$ & $<\brBTone$\\
$B_{T2}$ & $<\brBTtwo$\\
$E_0$ & $<\brEzero$\\
$E_1$ & $<\brEone$\\
$A_*$ & $<\brA$\\
$Y_0$ & $<\brYzero$\\
$Z$ & $<\brZ$\\
$M_2$ & $<\brMtwo$\\
$C_2$ & $<\brCtwo$\\
$Y_0+(Z-1)r_*+C_2r_*^2$ & $<\brPolynomial$\\
$Z+2C_2r_*$ & $<\brContraction$\\
$4C_2Y_0$ & $<\brFourCY$\\
\bottomrule
\end{tabular}
\caption{The certificate at $r_*=\brRadius$. The last row gives
$4C_2Y_0<(1-Z)^2$, the discriminant condition for the radii polynomial.}
\label{br:tab:constants}
\end{table}

\subsection{Geometry on the whole ball}\label{br:sec:geometry}
Let $\sigma_0,\sigma_1$ be the sums defined
after~\eqref{br:eq:complement-constants}.
For every shape in the ball, the $W_1$ variation norm bounds both
$\sum_j|p_j-p_j^{(0)}|$ and
$\sum_jjm|p_j-p_j^{(0)}|$ by $r_*$. Consequently, on $\overline\D$,
\begin{equation}\label{br:eq:whole-ball-geometry}
 |p|\ge\rho-\sigma_0-r_*>\brDerivativeLo,
 \qquad
 \re(1+wp'/p)\ge1-\frac{\sigma_1+r_*}{\rho-\sigma_0-r_*}
 >\brConvexLo.
\end{equation}
On the larger disc $|w|\le R$,
\begin{equation}\label{br:eq:univalence}
 \re p(w)\ge2\rho-P-r_*>\brUnivalentLo>0.
\end{equation}
For distinct $w_1,w_2$ inside that disc, the difference quotient
\[
 \frac{\psi(w_2)-\psi(w_1)}{w_2-w_1}
 =\int_0^1p(w_1+t(w_2-w_1))\,dt
\]
has positive real part, proving injectivity. Thus $\psi$ is univalent
and $p$ nonzero on a neighbourhood of $\overline\D$; $R>\brAnalyticRadius$.
The boundary is a real-analytic Jordan curve. The tangent-angle
calculation~\eqref{r2:eq:curvature-formula} is a purely conformal
identity and applies to the present boundary:
\begin{equation}\label{br:eq:curvature}
 \kappa(\theta)=\frac{\re(1+wp'/p)}{|p|}>0.
\end{equation}
Its positive curvature and total turning $2\pi$ make the bounded image
strictly convex.

For its first nonlinear coefficient,
\begin{equation}\label{br:eq:nondisc}
 |c_1|\ge|c_1^{(0)}|-\frac{r_*}{2(m+1)^2}>\brNonDiscLo.
\end{equation}
Indeed the $j=1$ term in $\nrm{p-p_0}_{W_1}$ is
$2(m+1)|p_1-p_1^{(0)}|$, and $c_1=p_1/(m+1)$.
The real coefficient and frequency restrictions give
\[
 \psi(e^{2\pi i/m}w)=e^{2\pi i/m}\psi(w),\qquad
 \psi(\bar w)=\overline{\psi(w)}.
\]
Thus $\Omega$ has $D_m$ symmetry. A disc invariant under this rotation
would be centred at zero, and its normalised Riemann map would be
linear, contradicting~\eqref{br:eq:nondisc}. Finally
$|\psi'(0)-\rho|\le r_*$ and the exact centre coefficient give
$\brRhoLo<\psi'(0)<\brRhoHi$.

\subsection{The field and its signs}\label{br:sec:sign}
At the zero, $U=\Bc(|p|)+Kg$ solves~\eqref{br:eq:pullback}
distributionally. Initially $U$ is continuous, with the Cauchy traces
given by the lift and Lemma~\ref{r2:lem:Poisson}. The coefficient
$|p|^2$ is real analytic near $\overline\D$. Dirichlet elliptic
regularity upgrades $U$ to a smooth solution on the closed disc,
and the Cauchy traces are classical. With
$u=U\circ\psi^{-1}$, conformal covariance and the normal transformation
give exactly~\eqref{br:eq:physical}.

At the origin the finite field has value $a_0/s_0$. Here $s_0$ is
the positive field divisor from~\eqref{br:eq:input}, while the shape
sums are $\sigma_0,\sigma_1$. The lift and ball bounds give
\begin{equation}\label{br:eq:origin}
 U(0)\ge\frac{a_0}{s_0}-L_0-r_*/4-B_bJ_br_*/2>\brOriginLo>0.
\end{equation}
For any boundary point $x$, smoothness and the outward derivative one
give $u(x-t\nu(x))=-t+O(t^2)<0$ for small positive $t$.
Thus $u$ changes sign. This completes the proof of
Theorem~\ref{br:thm:main} and, by the Green identity of
Subsection~\ref{br:sec:fourier}, its arc-length Fourier formulation.

\subsection{Relation with earlier work}\label{br:sec:literature}
\subsubsection{Sign, spectral sequences and low eigenvalues}
Serrin's symmetry theorem~\cite{Ser71}, in the form recorded
in~\cite[Theorem~2]{Dal14}, concerns a bounded connected open set
$\Omega$ with $C^2$ boundary and $u\in C^2(\overline\Omega)$ satisfying
$\Delta u+\lambda u+\varepsilon=0$, $\lambda>0$,
$\varepsilon\in\{0,\pm1\}$, and $u=0$, $\partial_\nu u=c$ on
$\partial\Omega$: if $u>0$ in $\Omega$, then $\Omega$ is a ball. Changing the sign of the solution gives the negative case.
The field in~\eqref{br:eq:origin} has both signs.
Berenstein's planar spectral-sequence theorem~\cite{Ber80}, as stated
in~\cite[Proposition~1.1]{Dal10} for simply connected bounded planar
domains with $C^{2,\varepsilon}$ boundary and in~\cite[\S0]{BY87}, forces a disc
when the Dirichlet overdetermined problem is solvable for infinitely
many distinct spectral parameters. Berenstein and Yang~\cite[Theorem 1]{BY87}
give the corresponding result for bounded Lipschitz domains with
connected boundary in all dimensions, for either the Neumann or
Dirichlet overdetermined problem. The present construction supplies
one parameter. These sequence results imply that the domain of
Theorem~\ref{br:thm:main} has only finitely many overdetermined Dirichlet
eigenvalues.

Dalmasso~\cite[Theorem~1]{Dal14} proves that a bounded connected domain
with $C^1$ boundary is a ball if $\Delta u+\lambda u=0$ has a
nonconstant solution $u\in C^2(\overline\Omega)$ with constant
Dirichlet and Neumann data and $\lambda\le\lambda_1$. If the domain
is also convex and symmetric about a hyperplane, the same conclusion
holds for $\lambda\le\lambda_2$.
The geometric and symmetry hypotheses of the latter statement hold
here. Its spectral inequality does not: the trial functions
$1-r^2$ and $r^n(1-r^2)\cos n\theta$, $r^n(1-r^2)\sin n\theta$,
$1\le n\le6$, span a thirteen-dimensional zero-Dirichlet space on
the disc. As in~\eqref{r2:eq:dirichlet-minmax}, their Rayleigh quotients
are $2(n+1)(n+3)\le126$. Conformal invariance of energy and
$|p|>\brDerivativeLo$ therefore give
\begin{equation}\label{br:eq:low-index}
 \lambda_2(\Omega)\le\lambda_{13}(\Omega)
 \le126/(\brDerivativeLo)^2<1.
\end{equation}
For the finite centre, comparison with the inscribed and circumscribed
origin-centred discs, using binary64 Bessel zeros, shows that between
$14037$ and $14040$ Dirichlet eigenvalues do not exceed one.
The second-Dirichlet-eigenvalue case of the Schiffer problem is recorded
in~\cite[\S0]{BY87} as a consequence of the Payne--Weinberger
isoperimetric inequality and attributed to Berenstein~\cite{Ber80}
in~\cite[Remark~1]{Dal14}. The corresponding application of Dai, Liu
and Sun~\cite{DLS26} is quoted in~\cite[\S1]{DSWZ25}. These statements
concern the Schiffer boundary data.

Dalmasso's finiteness theorem~\cite[Theorem 1.1]{Dal10} concerns
$\Delta u+\lambda u+\mu=0$ with $u=0$, $\partial_\nu u=c\ne0$.
For a bounded convex planar domain with smooth positive-curvature
boundary, the Schiffer property and nonconstant width, it allows at
most finitely many pairs $(\lambda,\mu)$ at fixed $c$.
It is a finiteness statement and permits the single pair $(1,0)$.
Liu's criterion~\cite{Liu07}, as recorded in~\cite{CSS26,DSWZ25},
involves constancy of the second interior normal derivative.
Here the boundary decomposition of the Laplacian gives directly
\begin{equation}\label{br:eq:second-normal}
 u_{\nu\nu}=-\kappa u_\nu=-\kappa\quad\text{on }\partial\Omega:
\end{equation}
the tangential second derivative of the zero boundary trace vanishes,
and $\Delta u=-u=0$ there. Constant curvature would make the boundary
a circle, so~\eqref{br:eq:nondisc} also excludes constancy of
$u_{\nu\nu}$.

\subsubsection{Dai--Sun--Wei--Zhang and the numerical remainder comparison}
Dai, Sun, Wei and Zhang~\cite[Theorem 1.1, version 2]{DSWZ25} prove
Dirichlet rigidity above a threshold $\alpha_0(\Omega)$ for bounded
uniformly convex planar domains with connected $C^{2,\varepsilon}$
boundary, $0<\varepsilon<1$. The geometric hypotheses hold here.
The threshold depends on the full geometry of the fixed domain.
In support coordinates in Steiner position, let $h(\theta)$ be the
support function and $\varrho=h+h''$ the radius of curvature.
Their amplitude defect is
\begin{equation}\label{br:eq:DSWZ-defect}
 \mathfrak B(\Omega)=\max_{\theta,\,\upsilon\in\{-1,1\}}
 \left|\sqrt{\varrho(\theta)}e^{\upsilon h'(\theta)}
 -\sqrt{\varrho(\theta+\pi)}e^{-\upsilon h'(\theta+\pi)}\right|.
\end{equation}
This is the quantity $\mathcal B(\Omega)$ of~\eqref{r2:eq:DSWZ-defect}.

For a controlled $C^5$ class, put
\[
 I_D(\lambda,\upsilon,\theta)=\int_{\partial\Omega}
 e^{\upsilon\eta_\perp\cdot x+i\lambda\xi\cdot x}\,ds,
 \qquad \xi=(\cos\theta,\sin\theta),\quad\eta_\perp\perp\xi.
\]
Its two-point leading term is
\begin{equation}\label{br:eq:DSWZ-leading}
 \sqrt{\frac{2\pi}{\lambda}}\left[
 \sqrt{\varrho(\theta)}e^{\upsilon h'(\theta)}e^{i\lambda h(\theta)-i\pi/4}
 +\sqrt{\varrho(\theta+\pi)}e^{-\upsilon h'(\theta+\pi)}
                          e^{-i\lambda h(\theta+\pi)+i\pi/4}\right].
\end{equation}
Writing $R_D$ for the difference, their expansion gives
$|R_D|\le C_*\lambda^{-3/2}$ for $\lambda\ge3$, $|\upsilon|\le1$.
The constant $C_*$ depends on the class and is not supplied numerically.
Their exclusion lemma, for $\mathfrak B\ge\delta>0$, has threshold
\begin{equation}\label{br:eq:DSWZ-threshold}
 \alpha>\left[\max\left\{3,\frac{C_*}{\sqrt{2\pi}\delta}\right\}\right]^2-1.
\end{equation}
The qualitative theorem chooses $\delta=\mathfrak B/2$.

The following comparison uses floating-point quadrature for the finite centre. Let $\Omega_1$ be its image rescaled to conformal radius
one; its parameter is $\alpha_u=\rho^2\approx\brAlphaDiagnostic$.
Any constant admissible for its stationary-phase expansion must
dominate every evaluated remainder. In the unit scale define
\[
 G_u(\zeta)=\zeta^{3/2}\max_\theta|R_D^{\Omega_1}(\zeta,0,\theta)|.
\]
The computed values in Table~\ref{br:tab:remainder} use trapezoidal
integration in the conformal angle and the support coordinates of the
two stationary points.
\begin{table}[htbp]
\centering
\begin{tabular}{rr}
\toprule
$\zeta$ & Sampled $G_u(\zeta)$\\
\midrule
\brRemainderRows
\bottomrule
\end{tabular}
\caption{Numerical stationary-phase remainders for the normalised
Berenstein centre in binary64 arithmetic. Each sampled angular maximum
estimates a lower bound for $G_u(\zeta)$.}
\label{br:tab:remainder}
\end{table}

The comparison also includes dilation. For $s\Omega_1$ the parameter
is $\alpha_u/s^2$, and
\[
 R_D^{s\Omega_1}(\lambda,0,\theta)
 =sR_D^{\Omega_1}(s\lambda,0,\theta),\qquad
 \mathfrak B(s\Omega_1)=\sqrt s\,\mathfrak B_u(s),
\]
where
$\mathfrak B_u(s)=\max_{\theta,\upsilon=\pm1}
|\sqrt\varrho\,e^{\upsilon sh'}-\sqrt{\varrho(\theta+\pi)}e^{-\upsilon sh'(\theta+\pi)}|$
uses the unit-scale support data. An admissible $C_*$ for a class
containing $s\Omega_1$ satisfies
\[
 C_*\ge\lambda'^{3/2}|R_D^{s\Omega_1}(\lambda',0,\theta)|,
 \qquad \lambda'\ge3.
\]
Taking the angular maximum gives
$C_*\ge s^{-1/2}G_u(s\lambda')$. With the most favourable choice
$\delta=\mathfrak B(s\Omega_1)$, exclusion of $\alpha_u/s^2$
by~\eqref{br:eq:DSWZ-threshold} requires
$\lambda_s=\sqrt{\alpha_u/s^2+1}>3$ and
\[
 C_*<\sqrt{2\pi}\lambda_s\mathfrak B(s\Omega_1)
 =\sqrt{2\pi(\alpha_u+s^2)}s^{-1/2}\mathfrak B_u(s).
\]
Hence exclusion requires
\begin{equation}\label{br:eq:DSWZ-rescaled}
 G_u(\zeta)<T_D(s):=\sqrt{2\pi(\alpha_u+s^2)}\,\mathfrak B_u(s)
 \quad\text{for every }\zeta\ge3s.
\end{equation}
For $s>\sqrt{\alpha_u/8}$ the parameter is below eight and
cannot exceed~\eqref{br:eq:DSWZ-threshold}. On the remaining scales
$0<s\le\sqrt{\alpha_u/8}\approx\brScaleDiagnostic$, the numerical
calculation gives $T_D(s)$ approximately in the range
$\brThresholdDiagnostic$. The sampled remainder at the largest
frequency in Table~\ref{br:tab:remainder} is approximately
$\brRemainderDiagnostic$, about $\brRatioDiagnostic$ times the
largest computed threshold. In this computation the condition
\eqref{br:eq:DSWZ-rescaled} fails at every admissible scale, so the
stationary-phase exclusion does not remove the finite centre.

Version~1 of~\cite{DSWZ25} required central symmetry. The present
domain is not centrally symmetric. Any centre of symmetry would be
fixed by its nontrivial rotational symmetry and hence be the origin.
Then normalised Riemann-map uniqueness would imply $\psi(-w)=-\psi(w)$.
Since $m=\brM$ is odd, its nonzero coefficient $c_1$ multiplies the
even power $w^{m+1}$, which contradicts that identity.

\subsubsection{The planar counterexample of Colbrook--Sadeghi--Stepaniants}
Colbrook, Sadeghi and Stepaniants~\cite{CSS26} proved existence on
a bounded simply connected non-disc domain with real-analytic Jordan
boundary, $D_{13}$ symmetry and a sign-changing eigenfunction.
For their numerical centre
$\varphi^\circ(w)=w+\sum_{j=1}^{20}\varphi_j^\circ w^{13j+1}$,
the binary64 evaluation of its appendix coefficients in the present paper gives
\begin{equation}\label{br:eq:CSS-computation}
 \min_{|w|=1}\re(1+w(\varphi^\circ)''/(\varphi^\circ)')
       \approx\brCSSConvexDiagnostic,\qquad
 \min\kappa^\circ\approx\brCSSCurvatureDiagnostic.
\end{equation}
Here $\kappa^\circ$ is the curvature of
$\varphi^\circ(\partial\D)$ in the normalisation
$(\varphi^\circ)'(0)=1$. The centre curve is not convex: one of its
points lies at distance about $1.0\cdot10^{-2}$ inside the boundary of
its convex hull. The exact boundary of~\cite{CSS26} differs from the
centre curve pointwise by less than $4.67\cdot10^{-10}$~\cite[Subsection~5.1]{CSS26},
so their domain is not convex. The interval triangle witness is given in
\path{anc/diagnostics/berenstein/csb_nonconvex.json}; its construction is
described in Subsection~\ref{br:sec:verification}.
Theorem~\ref{br:thm:main} supplies a convex domain.
The conformal fixed-disc method is due to~\cite{CS26},
and~\cite{CSS26} adapted it to the Dirichlet problem. The coupled inverse here uses the Dirichlet trace of
a Robin solution to recover the shape, together with the explicit
residual Neumann lift~\eqref{br:eq:corrected}.
